\section{Planteamiento del problema: el modelo no cambia, el presupuesto de cómputo aumenta}

Las scaling laws tradicionales destinan la mayor parte del cómputo a la etapa de entrenamiento; el test-time compute scaling, en cambio, pregunta: \textbf{una vez fijados los parámetros del modelo, ¿se puede elevar la tasa final de éxito del mismo modelo mediante más muestreo, búsqueda, revisión y verificación?} Esta sesión se organiza alrededor de tres líneas de trabajo:

\begin{enumerate}
    \item Large Language Monkeys: por qué el muestreo repetido presenta una scaling law predecible;
    \item Compute-optimal test-time scaling: cómo se debe distribuir el presupuesto entre el muestreo paralelo y la revisión secuencial;
    \item Archon: cómo buscar de manera automática una arquitectura de razonamiento compuesta por generator, critic, ranker, fuser y verifier.
\end{enumerate}

\begin{importantbox}{Descomposición central de esta sesión}
El escalamiento en tiempo de prueba no es un solo algoritmo, sino tres problemas: si la \textbf{generación} puede cubrir el candidato correcto, si la \textbf{verificación} puede identificar el candidato correcto y si el \textbf{control de presupuesto} puede destinar el cómputo limitado a las operaciones de mayor valor. Optimizar solo uno de estos tres aspectos puede hacer que los otros dos absorban la ganancia total.
\end{importantbox}

\subsection{Cuatro métricas de evaluación}

\begin{table}[H]
\centering
\begin{tabular}{p{3.1cm}p{10.5cm}}
\toprule
\textbf{Métrica} & \textbf{Significado} \\
\midrule
pass@1 & Precisión final cuando el sistema entrega una sola respuesta \\
pass@$k$ / coverage & Probabilidad de que, entre $k$ candidatos, exista al menos una respuesta correcta \\
best-of-$N$ & Después de generar $N$ candidatos, un verifier elige una respuesta \\
majority@$N$ & Se vota según la frecuencia de las respuestas, bajo el supuesto implícito de que «la respuesta correcta es la más frecuente» \\
\bottomrule
\end{tabular}
\end{table}

\subsection{Resumen de la sección}

Aumentar el cómputo de inferencia primero amplía el conjunto de candidatos, pero el sistema final entrega un solo resultado, por lo que la coverage y el pass@1 deben medirse por separado. Todos los métodos que siguen pueden entenderse como una manera de equilibrar la calidad de los candidatos, la calidad de la selección y el costo de cómputo.

